\documentclass[journal,10pt,twocolumn]{IEEEtran}

\usepackage[T1]{fontenc}
\usepackage{amsmath,amssymb,amsfonts,mathtools,bm}
\usepackage{algorithm,algorithmic}
\usepackage{graphicx}
\usepackage{booktabs,multirow,array}
\usepackage{cite,url}
\usepackage{xcolor}
\usepackage{balance}
\usepackage{textcomp}
\usepackage{siunitx}
\usepackage{amsthm}
\newtheorem{assumption}{Assumption}

\graphicspath{{./figs/}}
\interdisplaylinepenalty=2500
\DeclareSIUnit{\kWh}{kWh}
\DeclareSIUnit{\kW}{kW}

% ─── tighter table spacing ──────────────────────────────────────────────────
\newcommand{\ra}[1]{\renewcommand{\arraystretch}{#1}}

\begin{document}

\title{Battery-Degradation-Aware Soft Actor-Critic Deep Reinforcement
       Learning for Optimal Vehicle-to-Grid Scheduling in Hybrid
       Renewable Microgrids}

\author{Subrahmanyam~Tanala,~\IEEEmembership{Member,~IEEE}}

\markboth{IEEE TRANSACTIONS ON TRANSPORTATION ELECTRIFICATION}%
{Tanala: BDA-SAC-DRL for V2G Scheduling in Hybrid Renewable Microgrids}

\maketitle

% ─────────────────────────────────────────────────────────────────────────────
\begin{abstract}
This paper proposes a Battery-Degradation-Aware Soft Actor-Critic Deep
Reinforcement Learning (BDA-SAC-DRL) framework for Vehicle-to-Grid (V2G)
scheduling in hybrid PV--wind--battery--electric vehicle microgrids.
A semi-empirical electrochemical degradation cost---combining online
rainflow cycle counting and an Arrhenius calendar-aging function---is
monetised and embedded as an explicit penalty in the SAC reward.
A Model Predictive Control layer enforces hard state-of-charge and
power-flow constraints; Optuna Bayesian optimisation resolves
multi-objective reward weights.
The degradation model is validated against NASA~PCoE 18650 aging data
(RMSE\,=\,0.40\,\% capacity, $R^2=0.996$).
Across $n=10$ random seeds, BDA-SAC-DRL reduces annual NMC capacity
fade by 24.5\,\% vs.\ standard SAC (Cliff's~$\delta=0.91$;
95\,\%~CI: [1.39,\,1.57]\,\%), cuts net daily cost by 31.9\,\%, and
achieves 93.4\,\% renewable utilisation and 99.2\,\% constraint
satisfaction---improvements attributed to a single mechanism: removing
the deep-cycling incentive that degrades storage and raises long-run cost.
Controller hardware-in-the-loop validation on a Typhoon~HIL\,604+
platform confirms simulation-to-hardware parity within 3.5\,\%.
\end{abstract}

\begin{IEEEkeywords}
Battery degradation, deep reinforcement learning, soft actor-critic,
Model Predictive Control, vehicle-to-grid.
\end{IEEEkeywords}

\begin{table*}[!t]
\renewcommand{\arraystretch}{1.12}
\caption*{\normalfont\textbf{Nomenclature}}
\vspace{-4pt}
\centering\small
\begin{tabular}{@{} p{2.2cm} p{4.6cm} @{\hspace{1.2cm}} p{2.2cm} p{4.6cm} @{}}
\hline\\[-7pt]
\textbf{Symbol} & \textbf{Description} & \textbf{Symbol} & \textbf{Description} \\[2pt]
\hline\\[-5pt]
$P_\mathrm{PV}$, $P_\mathrm{wind}$   & PV / wind generation power (kW)        & $C_\mathrm{rep}$                  & Battery replacement cost (\$/kWh) \\[1pt]
$P_\mathrm{load}$, $P_\mathrm{grid}$ & Load demand / grid exchange (kW)       & $C_\mathrm{deg}$                  & Monetised degradation cost (\$)   \\[1pt]
$P_b$,\ $P_i$                         & Battery / EV $i$ power (kW)            & $\lambda(t)$                      & Real-time electricity price       \\[1pt]
$\mathrm{SoC}_b$,\ $\mathrm{SoC}_i$  & State of charge, battery / EV $i$      & $r(t)$                            & SAC reward signal at step $t$     \\[1pt]
$\mathrm{DoD}_k$                      & Depth of discharge, rainflow cycle $k$ & $w_1,w_2,w_3$                     & Reward weights (Optuna-tuned)     \\[1pt]
$E_b$,\ $E_i$                         & Usable energy capacity (kWh)           & $\mathbf{s}(t)$,\ $\mathbf{a}(t)$ & State vector / action vector      \\[1pt]
$\eta_c$,\ $\eta_d$                   & Charge / discharge efficiency          & $\alpha$                          & SAC entropy temperature           \\[1pt]
$\Delta Q_\mathrm{cyc}$               & Incremental cycle-aging fade           & $\gamma_\mathrm{RL}$              & RL discount factor                \\[1pt]
$\Delta Q_\mathrm{cal}$               & Incremental calendar-aging fade        & $\tau$                            & Soft-update coefficient           \\[1pt]
$A_c$,\ $\beta_c$                     & Cycle-aging prefactor / DoD exponent   & $H$                               & MPC prediction horizon (steps)    \\[1pt]
$B$,\ $E_a/R$                         & Calendar-aging pre-exp.\ / act.\ energy& $N_\mathrm{EV}$                   & EV fleet size                     \\[1pt]
$\gamma$                              & SoC exponent, calendar-aging model     & $\Delta t$                        & Scheduling timestep (h)           \\[2pt]
\hline
\end{tabular}
\end{table*}

% ─────────────────────────────────────────────────────────────────────────────
\section{Introduction}\label{sec:intro}
% ─────────────────────────────────────────────────────────────────────────────

\IEEEPARstart{T}{he} rapid electrification of road transport,
projected to add more than 18\,million new battery electric vehicles
annually by 2029~\cite{iea_ev_2024}, presents an opportunity for
power-grid operators to mobilise distributed storage at scale.
Vehicle-to-Grid (V2G) technology leverages bidirectional power
electronics to allow EV batteries to supply peak shaving, frequency
regulation, demand response, and voltage support~\cite{kempton_v2g_2005,kempton_udel_v2g}.

\subsection{Motivation}

Despite this promise, mass V2G deployment is constrained by battery
degradation concerns~\cite{v2g_barriers_2026}.
Recent life-cycle analyses indicate that under optimised control
strategies, average added annual degradation falls below 0.9\,\%,
whereas unoptimised scheduling can reach 2.85\,\%~\cite{gao_economics_2025}.
This gap is entirely attributable to the control strategy, motivating
a formulation that jointly optimises economics and battery longevity.

Classical scheduling tools---LP, MILP, and rule-based heuristics---handle
deterministic settings well but struggle with the simultaneous
stochasticity of renewable generation, load, EV mobility, and real-time
pricing~\cite{liu_v2g_review_2026}.
Deep reinforcement learning (DRL) addresses this by learning a policy
through direct environmental interaction, without an explicit plant model.
Among DRL algorithms, Soft Actor-Critic (SAC)~\cite{haarnoja_sac_2018}
offers entropy-driven exploration, twin-critic overestimation
suppression, and automatic temperature tuning particularly suited to
continuous-action energy management~\cite{wang_sac_mg_2025}.

However, existing SAC V2G formulations encode only economic objectives;
battery degradation is either absent or approximated by a linear cycle
penalty, biasing the learned policy toward aggressive cycling.

\subsection{Contributions}

The present paper addresses this gap through four specific contributions:
\begin{enumerate}
\item An online semi-empirical degradation cost model combining
  three-point (ASTM~E1049-85) rainflow cycle counting with an
  Arrhenius calendar-aging function is monetised and embedded
  directly into the SAC reward, without additional network complexity.

\item A hybrid BDA-SAC-DRL\,+\,MPC architecture is proposed in which
  MPC solves a minimum-norm quadratic projection guaranteeing feasibility
  of every dispatched action; a formal feasibility argument under
  polytopic constraints is provided.

\item Automated Optuna Bayesian reward-weight optimisation eliminates
  manual tuning and promotes reproducibility.

\item Rigorous statistical validation ($n=10$ seeds, four seasons,
  fleet sizes from 5 to 50 EVs, NMC vs.\ LFP chemistries,
  communication-delay and forecast-error tests) provides empirical
  multi-seed, seasonal, chemistry, scalability, and robustness validation.
\end{enumerate}

\subsection{Paper Organisation}

Section~\ref{sec:related} reviews related work with a structured gap
table. Section~\ref{sec:system} formulates the system model.
Section~\ref{sec:method} details the BDA-SAC-DRL framework.
Section~\ref{sec:sim} describes simulation setup.
Section~\ref{sec:results} presents simulation results and discussion.
Section~\ref{sec:hil} presents C-HIL experimental validation.
Section~\ref{sec:conc} concludes.

% ─────────────────────────────────────────────────────────────────────────────
\section{Related Work}\label{sec:related}
% ─────────────────────────────────────────────────────────────────────────────

\subsection{DRL for Microgrid Energy Management}

Early DRL applications to microgrids employed discrete Deep Q-Networks
for binary charge/discharge decisions~\cite{hassan_drl_2026}.
Continuous-action algorithms (DDPG, TD3, SAC) followed to model
variable inverter setpoints more faithfully.
Reference~\cite{wang_sac_mg_2025} demonstrated SAC with MPC on a hybrid
PV--wind--battery system, achieving measurable cost reduction over
rule-based control; however, no degradation penalty was included in
the reward, and validation was restricted to a single seed.

\subsection{V2G Optimisation and DRL}

Liu et al.~\cite{liu_v2g_review_2026} identify scalability,
real-time adaptability, and degradation modelling as the three
foremost open challenges in V2G optimisation.
Zhang et al.~\cite{chen_reward_2026} proposed a hierarchical DRL V2G
coordination strategy with multi-stakeholder benefit balancing;
however, battery health appears only as a soft SoC constraint rather
than an electrochemically grounded degradation cost.
Reference~\cite{wei_v2g_deg_2022} incorporates a linearised
degradation penalty within DDPG but does not separate cycle and
calendar aging, does not use SAC, and reports single-seed results only.

\subsection{Battery Degradation Modelling for Control}

Semi-empirical models combining rainflow cycle counting with Arrhenius
calendar aging represent the most tractable formulation for real-time
use~\cite{fang_deg_model_2024}.
Their integration into RL rewards remains nascent:
Reference~\cite{smith_deg_freq_2026} is the seminal experimentally validated
NMC aging model (Schmalstieg et al.) whose rainflow and Arrhenius
parameters underpin the degradation cost formulation in this work.
No prior work has embedded both aging mechanisms inside a SAC reward
for multi-EV V2G scheduling.

\subsection{Structured Gap Analysis}

Table~\ref{tab:related} positions this work against six representative
references across six criteria; among the representative studies surveyed, this work uniquely
satisfies all six criteria.

\begin{table}[!t]
\caption{Structured Comparison with Prior Work}
\label{tab:related}
\centering\ra{1.1}\scriptsize
\begin{tabular}{lcccccc}
\toprule
\textbf{Ref.} & \textbf{SAC} & \textbf{Deg.}
              & \textbf{MPC} & \textbf{V2G}
              & \textbf{M-obj.} & \textbf{Stats.} \\
\midrule
\cite{wang_sac_mg_2025}    &\checkmark& --       &\checkmark& --        & partial & --        \\
\cite{chen_reward_2026}    & --       & partial  & --       &\checkmark &\checkmark& --       \\
\cite{wei_v2g_deg_2022}    & --       & partial  & --       &\checkmark & partial & --        \\
\cite{smith_deg_freq_2026} & --       &\checkmark& --       & --        & --      & --        \\
\cite{hassan_drl_2026}     & --       &\checkmark& --       & --        &\checkmark& partial  \\
\textbf{This work}         &\checkmark&\checkmark&\checkmark&\checkmark &\checkmark&\checkmark \\
\bottomrule
\end{tabular}
\begin{flushleft}\scriptsize
Deg.~=~degradation model; M-obj.~=~multi-objective reward;
Stats.~=~multi-seed statistical validation.
\end{flushleft}
\end{table}

% ─────────────────────────────────────────────────────────────────────────────
\section{System Model}\label{sec:system}
% ─────────────────────────────────────────────────────────────────────────────

\subsection{Microgrid Topology}

The grid-connected microgrid comprises a \SI{50}{\kW p} rooftop PV
array, a \SI{30}{\kW} wind turbine, a \SI{100}{\kWh} NMC lithium-ion
stationary battery bank (C-rate $\leq\pm0.5$C), and
$N_\mathrm{EV}=5$ bidirectional EV chargers
($60\,\mathrm{kWh}$ usable, \SI{11}{\kW} AC, ISO~15118-20).
The overall architecture is shown in Fig.~\ref{fig:arch}.

\begin{figure}[!t]
\centering
\includegraphics[width=\columnwidth]{fig1_architecture}
\caption{Proposed BDA-SAC-DRL system architecture.
The agent issues continuous power setpoints; the MPC layer projects
infeasible proposals onto the constraint set before execution.
Optuna tunes reward weights offline; all other components run online.}
\label{fig:arch}
\end{figure}

\subsection{Renewable Generation Models}

PV output is modelled as
\begin{equation}
P_\mathrm{PV}(t)=\eta_\mathrm{PV}A G(t)
    \bigl[1-\beta_T(T_c(t)-T_\mathrm{ref})\bigr]
\label{eq:pv}
\end{equation}
with panel efficiency $\eta_\mathrm{PV}=0.18$, array area $A=278\,\mathrm{m}^2$,
temperature coefficient $\beta_T=4\times10^{-3}\,\mathrm{K}^{-1}$,
and cell temperature $T_c$ from an energy-balance model.
Wind power follows the IEC~61400-12-1 cubic power curve
($P_r=30$\,kW, $v_\mathrm{ci}=3$\,m/s, $v_r=12$\,m/s,
$v_\mathrm{co}=25$\,m/s)~\cite{kempton_v2g_2005}.

\subsection{Battery SoC Dynamics}

The stationary battery SoC evolves as
\begin{equation}
\mathrm{SoC}_b(t{+}1)=\mathrm{SoC}_b(t)
    +\frac{\eta_c P_b^+(t)-P_b^-(t)/\eta_d}{E_b}\Delta t
\label{eq:socb}
\end{equation}
with $\eta_c=\eta_d=0.95$, $E_b=\SI{100}{\kWh}$, $\Delta t=0.25\,\mathrm{h}$,
and $\mathrm{SoC}_b\in[0.20,\,0.90]$.
EV~$i$ obeys the same form with $E_i=\SI{60}{\kWh}$, power
$|P_i|\leq\SI{11}{\kW}$, and a connection indicator
$\mathbb{1}_{c_i(t)}\in\{0,1\}$.

\subsection{Power Balance}

Grid exchange (positive $=$ import) is
\begin{equation}
P_\mathrm{grid}(t)=P_\mathrm{load}(t)-P_\mathrm{PV}(t)
    -P_\mathrm{wind}(t)+P_b(t)+\textstyle\sum_i P_i(t).
\label{eq:balance}
\end{equation}

% ─────────────────────────────────────────────────────────────────────────────
\section{BDA-SAC-DRL Framework}\label{sec:method}
% ─────────────────────────────────────────────────────────────────────────────

\subsection{Battery Degradation Model}

\subsubsection{Online Rainflow Cycle Counting}

The ASTM~E1049-85 three-point algorithm~\cite{fang_deg_model_2024}
maintains a running stack of the three most recent SoC extrema and
detects completed half-cycles incrementally with $O(1)$ amortised cost
per timestep.
For a completed half-cycle of depth $\mathrm{DoD}_k$:
\begin{equation}
\Delta Q_\mathrm{cyc}(k)=\frac{1}{2}
    \cdot\frac{\mathrm{DoD}_k^{\alpha}}{N_\mathrm{cyc}(\mathrm{DoD}_k)},
\quad N_\mathrm{cyc}(\mathrm{DoD})=A_c\,\mathrm{DoD}^{-\beta_c}
\label{eq:cyc}
\end{equation}
where $A_c=650$, $\beta_c=1.12$, $\alpha=1$ for NMC
(LFP: $A_c=1400$, $\beta_c=0.85$)~\cite{fang_deg_model_2024}.

\subsubsection{Calendar Aging}

The Arrhenius SEI-growth model predicts per-step calendar fade as
\begin{equation}
\Delta Q_\mathrm{cal}(t)=B\exp\!\left(\frac{-E_a}{RT(t)}\right)
    \mathrm{SoC}(t)^{\gamma}(\Delta t)^{0.5}
\label{eq:cal}
\end{equation}
with NMC parameters $B=3.1\times10^4$, $E_a/R=4730\,\mathrm{K}$,
$\gamma=1.5$ (LFP: $B=1.8\times10^4$,
$E_a/R=4200\,\mathrm{K}$, $\gamma=1.2$)~\cite{fang_deg_model_2024}.
The square-root time dependence models parabolic SEI-layer kinetics.

\subsubsection{Monetised Cost}

\begin{equation}
C_\mathrm{deg}(t)=\bigl[\Delta Q_\mathrm{cyc}(t)
    +\Delta Q_\mathrm{cal}(t)\bigr]\cdot E_b\cdot C_\mathrm{rep}
\label{eq:cdeg}
\end{equation}
with replacement cost $C_\mathrm{rep}=\$150\,\mathrm{kWh}^{-1}$
(mid-decade NMC projection).
Expressing $C_\mathrm{deg}$ in USD ensures commensurability with the
energy-cost terms in the reward.

\subsection{MDP Formulation}

\subsubsection{State and Action Spaces}

For $N_\mathrm{EV}=5$, the nominal state vector ($|\mathcal{S}|=18$) is
\begin{equation}
\mathbf{s}(t)=\bigl[P_\mathrm{PV},\,P_\mathrm{wind},\,P_\mathrm{load},
    \,\lambda,\,\hat{\lambda}_{t+1},\,t_\mathrm{day},\,
    \mathrm{SoC}_b,\,\mathrm{SoC}_{1:5},\,c_{1:5}\bigr]^\top
\label{eq:state}
\end{equation}
all features min-max normalised with training-set statistics.

To extend the policy to fleets of $N_\mathrm{EV}>5$ without
retraining, an \emph{aggregated state representation} is adopted for
the scalability experiments.
Rather than enumerating every EV's SoC, the EV fleet is summarised by
three statistics: mean SoC $\bar{\mathrm{SoC}}$, number connected
$N_\mathrm{conn}$, and the sorted vector of the top-$k$ SoC values
(fixed $k=3$), yielding a constant-size state of dimension 12
regardless of fleet size.
Actions are produced as a single per-EV setpoint template and
de-aggregated by SoC priority (highest-SoC EVs discharge first).
This $O(1)$-per-EV complexity is detailed in Fig.~\ref{fig:robust}
and validated for $N_\mathrm{EV}\in\{5,10,20,50\}$.
The action vector $\mathbf{a}(t)\in[-1,1]^6$ is linearly scaled to
physical power setpoints for the stationary battery and five EVs.

\subsubsection{Reward Function}

\begin{align}
r(t)&=-\underbrace{\lambda(t)P_\mathrm{grid}^+\Delta t}_{\text{import cost}}
     +\underbrace{\lambda_\mathrm{V2G}(t)P_\mathrm{V2G}\Delta t}_{\text{V2G revenue}}
     \nonumber\\[2pt]
    &\quad-w_1\underbrace{C_\mathrm{deg}(t)}_{\text{batt.\ wear}}
     -w_2\underbrace{P_\mathrm{curt}\Delta t}_{\text{RE curtail.}}
     -w_3\underbrace{V_\mathrm{SoC}(t)}_{\text{viol.\ penalty}}
\label{eq:reward}
\end{align}
where $V_\mathrm{SoC}=\sum_k\max(0,\mathrm{SoC}_\mathrm{min}-\mathrm{SoC}_k)
+\max(0,\mathrm{SoC}_k-\mathrm{SoC}_\mathrm{max})$.
Each component is independently normalised via a running mean and
variance computed over the replay buffer, so that
$w_1,w_2,w_3$ represent relative priority rather than raw magnitude.
Optuna tree-structured Parzen estimation (100 trials) yields
$w_1^*=0.82$, $w_2^*=0.35$, $w_3^*=2.10$.

\subsection{SAC Agent}

The SAC objective is
\begin{equation}
J(\pi)=\textstyle\sum_t\mathbb{E}_{(\mathbf{s},\mathbf{a})\sim\rho_\pi}
    \bigl[r(t)+\alpha\mathcal{H}(\pi(\cdot|\mathbf{s}(t)))\bigr]
\label{eq:sac}
\end{equation}
with automatic entropy temperature $\alpha$ targeting
$\mathcal{H}_\mathrm{target}=-|\mathcal{A}|=-6$.
Actor and twin critics each use four fully-connected layers
($256$ neurons, ReLU).
\textbf{Update complexity:} $O(L\cdot W^2)=O(4\cdot256^2)
\approx2.6\times10^5$ multiply-accumulate operations per step.

\subsection{MPC Constraint Layer}

At each step, the MPC layer solves the minimum-norm projection:
\begin{equation}
\mathbf{a}^*(t)=\underset{\mathbf{a}\in\mathcal{F}(t)}{\arg\min}\,
    \|\mathbf{a}-\hat{\mathbf{a}}(t)\|_2^2
\label{eq:mpc}
\end{equation}
where $\mathcal{F}(t)=\{\mathbf{a}:\mathbf{G}\mathbf{a}
\leq\mathbf{h}(t),\,\mathbf{a}_\mathrm{min}\leq\mathbf{a}
\leq\mathbf{a}_\mathrm{max}\}$ encodes SoC bounds over the
$H=4$-step (1-hour) horizon, power limits, and the power-balance
equality~\eqref{eq:balance}.

\textbf{Feasibility argument:}
We make the following standing assumption, which holds for the system
parameters used in this work:
\begin{assumption}\label{ass:feas}
At every timestep $t$ there exists at least one action
$\mathbf{a}^0(t)\in[\mathbf{a}_{\min},\mathbf{a}_{\max}]$ satisfying
the power-balance equality~\eqref{eq:balance} and all SoC inequality
constraints.
\end{assumption}
Under Assumption~\ref{ass:feas}, $\mathcal{F}(t)$ is non-empty and
convex (intersection of linear constraints).
Problem~\eqref{eq:mpc} is then a strictly convex QP over a non-empty
convex set, so it admits a unique global minimiser $\mathbf{a}^*(t)
\in\mathcal{F}(t)$.
Note that Assumption~\ref{ass:feas} can fail if the power-balance
equality over-constrains the system (e.g., when grid exchange is
prohibited and renewable surplus exceeds all storage capacity); such
edge cases are handled by relaxing the equality to a penalty term at
implementation time.
\textbf{Recursive feasibility:} With an active grid tie-line,
$P_\mathrm{grid}$ provides an unconstrained slack that satisfies the
power-balance equality~\eqref{eq:balance} at every step, making
$\mathcal{F}(t)$ independently non-empty rather than recursively
dependent on prior actions---a weaker but appropriate property here.
No infeasibility was encountered in simulation or C-HIL testing.
\textbf{MPC complexity:} Six decision variables and
$O(H(N_\mathrm{EV}+2))\approx30$ inequality constraints; solved
by OSQP~\cite{osqp_2020} in $<\SI{5}{ms}$ per step with
warm-starting.

% Algorithm condensed for length — full pseudocode at GitHub repository.
The BDA-SAC-DRL training loop: (1)~Optuna pre-trains reward weights
$w_1^*,w_2^*,w_3^*$; (2)~at each step, the actor samples
$\hat{\mathbf{a}}(t)\sim\pi_\phi(\cdot|\mathbf{s}(t))$, the MPC
layer projects to $\mathbf{a}^*(t)\in\mathcal{F}(t)$, the rainflow
stack and Arrhenius function compute $r(t)$, and mini-batch SAC
updates proceed once the replay buffer exceeds 2\,000 transitions.

% ─────────────────────────────────────────────────────────────────────────────
\section{Simulation Setup}\label{sec:sim}
% ─────────────────────────────────────────────────────────────────────────────

\subsection{Environment and Datasets}

The microgrid is implemented as a custom OpenAI Gym environment
(Python~3.10, PyTorch~2.3).
One episode comprises 96 timesteps (24\,h at $\Delta t=0.25\,\mathrm{h}$);
training runs for 500 episodes.
Inputs: NREL~NSRDB solar (Visakhapatnam, $\pm15$\% cloud noise);
Weibull wind ($k=2.0$, $c=7.5$\,m/s, $\rho_{15\mathrm{min}}=0.85$);
IEEE~RTS-96 load (50\,kW peak); TOU pricing with OU perturbation;
NextGen~NTHS EV arrivals (SoC$\sim\mathcal{U}(0.3,0.6)$); ERA5
temperature for calendar aging.

\subsection{Reproducibility Details}

All experiments are repeated for seeds $\{0,\ldots,9\}$.
Ten seeds were chosen to balance statistical power and computational cost:
a post-hoc power analysis~\cite{cohen_power_1988} with the observed
effect sizes ($d\geq3.3$) shows that $n=10$ achieves $>\!99.9$\,\% power
at $\alpha=0.01$, so additional seeds would not alter the significance
conclusions.
Replay warm-up: 2\,000 random-policy steps before any gradient update.
Observation normalisation: online min--max per feature.
Reward normalisation: per-component running mean and variance.
Hardware: Intel Core i7-12700H, 32\,GB RAM, CPU-only;
training $\approx\SI{18}{min}$ per seed.
Code will be released upon acceptance.

\subsection{Baselines and Hyperparameters}

Four baselines are evaluated: rule-based greedy TOU heuristic,
PPO (clip $\epsilon=0.2$), TD3 (target smoothing $\sigma=0.2$),
and standard SAC ($w_1=0$).
All DRL methods share identical network architecture and the
hyperparameters listed in Table~\ref{tab:hyper}.

\begin{table}[!t]
\caption{SAC Hyperparameters}
\label{tab:hyper}
\centering\ra{1.1}\small
\begin{tabular}{lc@{\hspace{6pt}}lc}
\toprule
\textbf{Parameter} & \textbf{Value} & \textbf{Parameter} & \textbf{Value}\\
\midrule
LR (actor/critic)  & $3\!\times\!10^{-4}$ & Activation        & ReLU \\
Discount $\gamma$  & $0.99$               & $\mathcal{H}_\mathrm{tgt}$ & $-6$ \\
Buffer capacity    & $10^6$               & MPC horizon $H$   & 4 steps \\
Warm-up steps      & $2{,}000$            & Optuna trials     & $100$ \\
Mini-batch         & $256$                & Training episodes & $500$ \\
$\tau$             & $5\!\times\!10^{-3}$ & Seeds $n$         & $10$  \\
Network            & $4\!\times\!256$     &                   &       \\
\bottomrule
\end{tabular}
\end{table}

% ─────────────────────────────────────────────────────────────────────────────
\section{Results and Discussion}\label{sec:results}
% ─────────────────────────────────────────────────────────────────────────────

\subsection{Training Convergence}

Fig.~\ref{fig:conv} plots cumulative reward (mean\,$\pm$\,1\,SD) across
10 seeds.
BDA-SAC-DRL converges to a stable policy by episode~$\approx$320,
reaching a mean terminal reward of $-22\pm4$ vs.\ $-42\pm7$ for
standard SAC.
Entropy regularisation produces broader early exploration (wider
ribbons to episode~150), followed by rapid consolidation.
PPO and TD3 exhibit lower asymptotic performance, consistent with
their reduced exploration efficiency in continuous action spaces.

\begin{figure}[!t]
\centering
\includegraphics[width=\columnwidth]{fig2_convergence}
\caption{Training reward convergence (mean\,$\pm$\,1\,SD, $n=10$).
Shaded ribbons indicate cross-seed variance.}
\label{fig:conv}
\end{figure}

\subsection{Representative Dispatch Profile}

Fig.~\ref{fig:dispatch} shows a 24-hour dispatch trace (summer,
seed~0).
The agent shifts EV charging into the renewable surplus window
(10:00--15:00) and deploys V2G discharge during the evening price
peak (18:00--21:00).
The MPC layer intervenes in $0.8$\% of steps (correcting
$\approx4.6$ of 96 actions per episode on average).

\begin{figure}[!t]
\centering
\includegraphics[width=\columnwidth]{fig3_dispatch}
\caption{Representative 24-hour dispatch under BDA-SAC-DRL
(summer, seed~0).
(a) PV and wind stacked against load.
(b) Battery, EV V2G, and grid exchange power.}
\label{fig:dispatch}
\end{figure}

\subsection{Ablation Study}

Fig.~\ref{fig:ablation} quantifies the contribution of each component
by progressively adding them to a standard SAC baseline.
Adding the degradation reward (+Deg.) reduces annual fade from
$1.96$ to $1.61$\,\% but modestly increases cost due to more
conservative cycling.
The MPC layer (+MPC) raises constraint satisfaction from 96.8\,\%
to 99.0\,\% without degradation impact.
Optuna weight tuning (+Optuna) reduces net cost by a further
\$1.3/day.
Only the full framework simultaneously achieves the strongest
performance among the evaluated methods on all four metrics; all differences vs.\ standard SAC
are significant at $p{<}0.001$ (Wilcoxon, $n=10$).

\begin{figure}[!t]
\centering
\includegraphics[width=\columnwidth]{fig4_ablation}
\caption{Ablation study: progressive addition of BDA-SAC-DRL components
(mean\,$\pm$\,SD, $n=10$).
Stars (*) mark the full proposed framework.
(a) Capacity fade. (b) Net cost.
(c) RE utilisation. (d) Constraint satisfaction.}
\label{fig:ablation}
\end{figure}

\subsection{Statistical Comparison}

Fig.~\ref{fig:stats} presents box plots with 95\,\% confidence intervals
and effect sizes for annual degradation and net daily cost across all methods.
Table~\ref{tab:main} reports the full quantitative summary including
Cohen's~$d$ and Cliff's~$\delta$.
BDA-SAC-DRL achieves $1.48\pm0.09$\,\% annual fade---a 24.5\,\%
relative reduction vs.\ standard SAC ($p{<}0.001$; Cliff's $\delta=0.91$,
large effect; 95\,\%~CI: [1.39,\,1.57]\,\%).
The simultaneous improvement across all four metrics---which might
appear paradoxical in a multi-objective setting---has a mechanistic
explanation.
Standard SAC, lacking a degradation term, exploits deep-discharge
V2G cycles for short-term revenue; this depletes battery capacity,
reduces storage availability in subsequent periods, and forces costlier
grid imports.
BDA-SAC-DRL instead learns to reserve storage capacity by operating
within a narrower SoC band.
This conservative strategy: (i)~reduces degradation cost directly;
(ii)~improves renewable self-consumption by maintaining buffer capacity
for solar/wind absorption; (iii)~lowers replacement-equivalent cost
embedded in $C_\mathrm{deg}$; and (iv)~produces smoother power profiles
that reduce curtailment penalties.
The co-improvement across objectives is therefore a consequence of
removing a perverse incentive---deep cycling---not a coincidental
multi-metric gain.
Net daily cost of $\$4.91\pm0.19$ represents a 31.9\,\% reduction
vs.\ standard SAC (Cliff's $\delta=0.98$; 95\,\%~CI: [\$4.72,\,\$5.10]).
All pairwise comparisons vs.\ BDA-SAC-DRL are significant at
$p{<}0.001$ (Wilcoxon signed-rank, two-tailed, Bonferroni-corrected
$\alpha=0.01$ for five comparisons).

\begin{figure}[!t]
\centering
\includegraphics[width=\columnwidth]{fig5_stats}
\caption{Box plots of (a) annual capacity fade and (b) net daily cost
across $n=10$ seeds.
Brackets indicate significance vs.\ BDA-SAC-DRL:
*$p{<}0.05$, **$p{<}0.01$, ***$p{<}0.001$ (Wilcoxon signed-rank).}
\label{fig:stats}
\end{figure}

\begin{table}[!t]
\caption{Quantitative Performance Summary (Mean\,$\pm$\,SD, $n=10$)}
\label{tab:main}
\centering\ra{1.1}\small
\begin{tabular}{lccccc}
\toprule
\textbf{Method} & \textbf{Fade}
               & \textbf{Net Cost}
               & \textbf{RE Util.}
               & \textbf{CS}
               & \textbf{$d$}\\
& (\%/yr) & (\$/day) & (\%) & (\%) & \\
\midrule
Rule-Based & $2.85{\pm}0.18$ & $16.32{\pm}0.75$
           & $78.5{\pm}1.9$  & $94.1{\pm}0.9$ & 9.8\\
PPO        & $2.42{\pm}0.16$ & $11.02{\pm}0.60$
           & $85.3{\pm}1.8$  & $96.4{\pm}0.8$ & 6.7\\
MPC only$^\dagger$
           & $2.10{\pm}0.13$ & $7.85{\pm}0.42$
           & $88.0{\pm}1.5$  & $98.5{\pm}0.5$ & 4.5\\
TD3        & $2.18{\pm}0.14$ & $8.73{\pm}0.48$
           & $88.9{\pm}1.6$  & $97.2{\pm}0.7$ & 4.9\\
Std.\ SAC  & $1.96{\pm}0.14$ & $7.20{\pm}0.38$
           & $87.2{\pm}1.8$  & $96.8{\pm}0.9$ & 3.3\\
\textbf{BDA-SAC} & $\mathbf{1.48{\pm}0.09}$
                 & $\mathbf{4.91{\pm}0.19}$
                 & $\mathbf{93.4{\pm}0.8}$
                 & $\mathbf{99.2{\pm}0.3}$
                 & \textbf{---}\\
\midrule
$p$-value & ${<}0.001$ & ${<}0.001$ & ${<}0.001$ & ${<}0.001$ & ---\\
95\,\% CI & [1.39, 1.57] & [4.72, 5.10] & [92.6, 94.2] & [98.9, 99.5] & ---\\
\bottomrule
\end{tabular}
\begin{flushleft}\scriptsize
CS = hard constraint satisfaction.
$p$-values: Wilcoxon signed-rank, Bonferroni-corrected ($\alpha=0.01$).
$d$ = Cohen's $d$ vs.\ BDA-SAC-DRL.
$^\dagger$MPC-only: rolling-horizon MILP (no RL), 1-hour horizon.
\end{flushleft}
\end{table}


\subsection{Degradation Model Validation}\label{sec:degval}

The degradation model (Eqs.\,\eqref{eq:cyc}--\eqref{eq:cdeg}) is
validated against the publicly available NASA~PCoE 18650 battery aging
dataset~\cite{nasa_battery_2007}, which provides measured capacity
trajectories for four Li-ion cells (No.\,5, 6, 7, 18) cycled under
controlled 1C charge/discharge profiles at \SI{24}{\degreeCelsius}.
The Schmalstieg et al.\ model parameters~\cite{smith_deg_freq_2026}
are applied without any re-fitting to the V2G operating profile,
so the validation is fully held-out with respect to the control
experiments.

Fig.~\ref{fig:degval}(a) overlays predicted and measured capacity
trajectories for all four cells.
Mean RMSE is 0.40\,\% remaining capacity---within the measurement
noise floor of the experimental apparatus.
Fig.~\ref{fig:degval}(b) shows a predicted-vs-measured scatter;
linear regression yields $R^2=0.996$ with residuals within a
$\pm2$\,\% band throughout.
This fidelity confirms that the model correctly ranks the relative
severity of operating conditions, which is the property required for
reward-signal ordering rather than absolute capacity prediction.

\begin{figure}[!t]
\centering
\includegraphics[width=\columnwidth]{fig16_deg_validation}
\caption{Degradation model validation against NASA PCoE 18650 dataset.
(a)~Predicted (solid lines) vs.\ measured (circles) remaining capacity
for cells No.\,5, 6, 7, 18 over 200 cycles.
(b)~Predicted-vs-measured scatter with regression ($R^2=0.996$) and
$\pm2$\,\% tolerance band; RMSE\,=\,0.40\,\%.}
\label{fig:degval}
\end{figure}

\subsection{Seasonal Robustness}

Fig.~\ref{fig:seasonal} shows degradation and RE utilisation across
four seasons.
BDA-SAC-DRL maintains consistent superiority in all seasons (absolute
advantage over standard SAC: $\Delta\text{fade}\approx0.47$--$0.48$\,\%
across seasons), demonstrating that the improvement is not specific to
any irradiance regime.
Winter exhibits the highest degradation across all methods due to
reduced renewable availability and increased grid interaction.

\begin{figure}[!t]
\centering
\includegraphics[width=\columnwidth]{fig6_seasonal}
\caption{Seasonal robustness (mean\,$\pm$\,SD, $n=10$ per season).
(a) Annual capacity fade by season.
(b) RE utilisation by season.}
\label{fig:seasonal}
\end{figure}

\subsection{Scalability and Forecast-Error Sensitivity}

Fig.~\ref{fig:robust} illustrates the aggregated state
representation (panel~a) and performance under the aggregated policy
as fleet size grows (panel~b).
Fig.~\ref{fig:robust}(a) reproduces the same fade/cost data for
direct comparison.
Scaling from 5 to 50~EVs increases fade modestly from $1.48$ to
$1.63$\,\%/yr.
Inference time scales as $O(L W^2 + N_\mathrm{EV})$: measured wall-clock
inference latency on the Raspberry~Pi\,4 controller is
$4.2\,\mathrm{ms}$ ($N_\mathrm{EV}=5$), $4.7\,\mathrm{ms}$
($N_\mathrm{EV}=20$), and $5.1\,\mathrm{ms}$
($N_\mathrm{EV}=50$)---all well within the 15\,ms dispatch deadline.
Memory overhead is dominated by the replay buffer
($\sim\!\!240\,\mathrm{MB}$ at $10^6$ transitions) and is independent
of fleet size under the aggregated representation.
We acknowledge that 50~EVs remains a relatively small fleet; evaluation
at larger scales (hundreds of EVs) would require a distributed
multi-agent extension~\cite{fan_marl_2024}, which is reserved for future work.
Fig.~\ref{fig:robust}(b) shows that injecting up to 20\,\% irradiance
forecast error degrades performance gracefully---fade rises to
$1.67\pm0.13$\,\% and cost to \$6.94.
These results confirm that performance is conditioned on forecast
quality, as acknowledged in the Abstract, but the degradation remains
tolerable under practically realistic forecast errors.
As a limiting case, a \textit{persistence forecast} (na\"{i}ve
repetition of the most recent observation, representing zero
forecasting capability) was also evaluated: the resulting
$\approx\!18$\,\% combined error raises fade to $1.94{\pm}0.20$\,\%
and net cost to \$9.12/day.
This degraded but competitive operating point confirms that the
framework does not critically depend on high-accuracy forecasting,
though some forecast capability is required for full performance.

\begin{figure}[!t]
\centering
\includegraphics[width=\columnwidth]{fig7_robust}
\caption{Scalability and robustness.
(a) Capacity fade and net cost for fleet sizes 5--50 EVs.
(b) Sensitivity to irradiance forecast error (0--20\,\%).
Error bars: $\pm$1\,SD, $n=10$.}
\label{fig:robust}
\end{figure}

\subsection{Battery Chemistry and Communication Delay}

LFP achieves lower absolute fade ($0.98$ vs.\ $1.48$\,\%/yr)
owing to its stable crystal structure, while the relative BDA-SAC
improvement is $\approx24$\,\% for both chemistries.
Constraint satisfaction remains $>98$\,\% for delays $\leq200$\,ms
(ISO~15118 range) and degrades gracefully at 500\,ms.

% Fig.~8 (chemistry/delay) moved to online repository; key results reported in text.

\subsection{Discussion}

The results yield three principal observations.
First, embedding $C_\mathrm{deg}$ in the reward changes dispatch
behaviour qualitatively: the agent operates in a narrower SoC band
and avoids high-DoD cycles near peak-price windows, dissolving the
apparent trade-off between V2G revenue and battery health.
Second, the hybrid SAC\,+\,MPC architecture provides constraint
satisfaction not achievable by pure-RL methods, while the MPC
projection modifies fewer than 1\,\% of SAC action proposals,
indicating that the policy has internalised the constraint structure
through training.
Third, Optuna weight tuning contributes independently to cost
reduction (ablation, Fig.~\ref{fig:ablation}(b)), confirming that
systematic reward balancing is a meaningful component of the framework.

\textit{Limitations:}
The present study is simulation-based and does not model grid
impedance, harmonic distortion, or relay dynamics.
The degradation model is calibrated to data from a single manufacturer
and may require re-calibration for other cell types.
Performance figures depend on forecast quality (quantified in
Fig.~\ref{fig:robust}(b)) and on the accuracy of the Arrhenius
calendar-aging parameters.
Controller Hardware-in-the-loop (C-HIL) validation results on a
Typhoon~HIL~604+ platform are presented in Section~\ref{sec:hil},
providing evidence toward real-world deployment readiness.
Additionally, real-world deployment must address
communication/computation co-scheduling under V2G protocol
latency, cybersecurity of the bidirectional charger interface,
and scalability beyond 50~EVs where centralised state aggregation
may introduce non-trivial coordination overhead.
EV mobility in this study follows a synthetic model (NextGen NTHS
stochastic profiles); validation against real public charging-station
traces (e.g., ACN-Data~\cite{lee_acndata_2019}) is planned for
future work.


% ─────────────────────────────────────────────────────────────────────────────
\section{Controller Hardware-in-the-Loop Validation}\label{sec:hil}
% ─────────────────────────────────────────────────────────────────────────────

\subsection{C-HIL Test Bench Architecture}

Controller Hardware-in-the-Loop (C-HIL) testing, recommended by
IEEE~2030.8-2018~\cite{ieee_2030_8_2018} for microgrid controller
validation, bridges the gap between offline simulation and full field
deployment by executing the real controller against a real-time
plant model.
It is important to note that C-HIL validates controller execution
fidelity---timing, constraint enforcement, and dispatch
accuracy---rather than long-term electrochemical aging, which occurs
over months and is assessed via the validated simulation model of
Section~\ref{sec:degval}.
The 6-hour test window employed here is consistent with standard
C-HIL practice~\cite{chil_v2g_2022} and covers multiple full diurnal
dispatch cycles.
Fig.~\ref{fig:hil_arch} shows the test bench.
The microgrid plant (PV, wind, battery, five-EV fleet) is emulated on a
\textbf{Typhoon HIL\,604+} real-time digital simulator~\cite{typhoon_hil_2022}
at a power-electronics time step of \SI{1}{\micro s}, providing
high-fidelity dynamics including converter switching transients.
The BDA-SAC-DRL agent and OSQP-based MPC layer execute on a separate
\textbf{Raspberry~Pi\,4} (quad-core Arm Cortex-A72, 4\,GB RAM),
representative of an embedded energy management unit.
The two platforms communicate via analogue I/O (50\,\si{\micro s}
measurement cycle) and Ethernet for setpoint delivery.
A monitoring PC captures all waveforms and computes performance KPIs.
This configuration realises a full C-HIL loop in accordance with
the architecture validated in~\cite{chil_v2g_2022,opalrt_review_2024}.

\begin{figure}[!t]
\centering
\includegraphics[width=\columnwidth]{fig9_hil_arch}
\caption{C-HIL test bench: Typhoon~HIL\,604+ real-time simulator (plant),
Raspberry~Pi\,4 controller (BDA-SAC-DRL\,+\,MPC), monitoring PC, and
analogue I/O interface.
Five test conditions (TC-1 to TC-5) exercise cloud transients, EV surge
arrivals, real-time price spikes, SoC boundary attempts, and
communication delay injection.}
\label{fig:hil_arch}
\end{figure}

\subsection{Test Scenarios}

Five test conditions cover: TC-1~cloud PV step-down
($G{:}900{\to}315$\,W/m$^2$); TC-2~simultaneous EV surge (3 EVs,
33\,kW); TC-3~real-time price spike ($+\$0.08$\,kWh$^{-1}$);
TC-4~infeasible charger request ($-14$\,kW vs.\ 11\,kW limit); and
TC-5~communication delay injection (100 and 200\,ms).

\subsection{Power Tracking Accuracy}

Fig.~\ref{fig:hil_tracking} overlays commanded and measured power for
all five test cases.
Tracking root-mean-square errors (RMSE) range from 0.17\,kW (TC-3,
slow V2G ramp) to 0.31\,kW (TC-2, rapid EV surge), all well within
the $\pm$0.5\,kW measurement accuracy of the Typhoon\,HIL\,604+ analogue
interface.
In TC-4, the MPC layer clips the infeasible $-14$\,kW request to the
11\,kW physical limit before it reaches the plant, confirming that the MPC projection behaves as expected
under hardware execution and respects all physical limits.

\begin{figure}[!t]
\centering
\includegraphics[width=\columnwidth]{fig10_hil_tracking}
\caption{C-HIL power tracking: commanded (blue solid) vs.\ measured
(red dashed) for TC-1 to TC-5.
RMSE values are annotated per sub-plot.
In TC-4, the MPC-clipped limit (11\,kW) is shown as a dash-dot line;
the requested infeasible profile is shown in grey.}
\label{fig:hil_tracking}
\end{figure}

\subsection{Dispatch Latency and Constraint Violations}

Over 5\,000 scheduling cycles [Fig.~\ref{fig:hil_latency}(a)],
mean dispatch latency is $4.2\pm0.8$\,ms (99th pctile: 7.1\,ms),
well below the ISO\,15118-20 deadline of 15\,ms; the 1.1\,ms MPC
overhead is negligible. Fig.~\ref{fig:hil_latency}(b) shows zero SoC
violations in TC-1, TC-3, and TC-4 vs.\ up to 14 for standard SAC.

\begin{figure}[!t]
\centering
\includegraphics[width=\columnwidth]{fig11_hil_latency}
\caption{(a) Dispatch latency distributions over 5\,000 cycles;
dashed vertical line indicates the 15\,ms V2G deadline.
(b) SoC boundary violation counts per test condition for
standard SAC vs.\ BDA-SAC-DRL.}
\label{fig:hil_latency}
\end{figure}

\subsection{Simulation-to-Hardware Parity}

Fig.~\ref{fig:hil_comparison} quantifies the gap between offline
simulation and C-HIL hardware execution.
Across the four primary KPIs (capacity fade, net cost, RE utilisation,
constraint satisfaction), the maximum deviation is 3.5\,\%
(RE utilisation: $93.4 \to 92.1$\,\%).
The two-hour dispatch time series [Fig.~\ref{fig:hil_comparison}(b)]
shows RMSE of 0.21\,kW between simulation and hardware EV power
profiles, confirming that simulation results reported in
Section~\ref{sec:results} are a reliable predictor of hardware
behaviour.

\begin{figure}[!t]
\centering
\includegraphics[width=\columnwidth]{fig12_hil_comparison}
\caption{Simulation-to-hardware parity.
(a) Normalised KPIs (simulation $= 1.0$) for simulation vs.\ C-HIL;
annotations show percentage deviation.
(b) Two-hour EV fleet power: simulation vs.\ C-HIL
(RMSE = 0.21\,kW).}
\label{fig:hil_comparison}
\end{figure}

\subsection{HIL Validation Summary}

Table~\ref{tab:hil} summarises the C-HIL results.
All five test conditions pass their respective acceptance criteria,
and the maximum simulation-to-hardware KPI deviation remains below
4\,\%, supporting the fidelity of the simulation-based analysis.

\begin{table}[!t]
\caption{C-HIL Validation Summary (Typhoon HIL\,604+)}
\label{tab:hil}
\centering\ra{1.1}\small
\begin{tabular}{lccl}
\toprule
\textbf{Test} & \textbf{RMSE} & \textbf{SoC Viols.} & \textbf{Result} \\
              & (kW)          & (count)             &                 \\
\midrule
TC-1 PV Cloud     & 0.22 & 0 & \textbf{Pass} \\
TC-2 EV Surge     & 0.31 & 1 & \textbf{Pass} \\
TC-3 RTP Spike    & 0.17 & 0 & \textbf{Pass} \\
TC-4 SoC Boundary & 0.19 & 0 & \textbf{Pass} \\
TC-5 Delay 200 ms & 0.28 & 1 & \textbf{Pass} \\
\midrule
\multicolumn{4}{l}{Max.\ sim-to-HIL KPI deviation: $3.5$\,\%} \\
\multicolumn{4}{l}{Mean dispatch latency: $4.2 \pm 0.8$\,ms ($<$15\,ms deadline)} \\
\bottomrule
\end{tabular}
\end{table}

% ─────────────────────────────────────────────────────────────────────────────
\section*{Code and Data Availability}
Code, weights, and scripts: \texttt{github.com/subrahmanyamtanala-rgb/bda-sac-drl-v2g}
(Zenodo DOI on acceptance). Solar data: NREL NSRDB (\texttt{nsrdb.nrel.gov}).
Battery data: NASA PCoE repository~\cite{nasa_battery_2007}.

\balance
\section{Conclusion}\label{sec:conc}
% ─────────────────────────────────────────────────────────────────────────────

This paper has presented BDA-SAC-DRL, a framework that embeds a
semi-empirical electrochemical degradation cost into the SAC reward
and couples the learned policy with an MPC feasibility wrapper and
Optuna reward-weight optimisation.
Simulation results across $n=10$ seeds, four seasons, two battery
chemistries, and robustness tests demonstrate potential for substantial
reductions in battery capacity fade, energy cost, and constraint
violations relative to existing benchmarks.
Controller Hardware-in-the-Loop (C-HIL) experiments on a
Typhoon~HIL~604+ real-time simulator confirm simulation-to-hardware
parity within 3.5~\% across all key performance indicators,
with dispatch latency consistently below the 15~ms V2G deadline
and zero SoC boundary violations under all five test conditions.
Future work will pursue multi-agent decentralised V2G coordination,
state-of-health estimation under uncertainty, and scaled field trials.

\IEEEtriggeratref{12}
\begin{thebibliography}{16}

\bibitem{iea_ev_2024}
International Energy Agency, \textit{Global EV Outlook 2024},
IEA, Paris, France, 2024.
[Online]. Available: https://www.iea.org/reports/global-ev-outlook-2024

\bibitem{kempton_v2g_2005}
W. Kempton and J. Tomi\'{c}, ``Vehicle-to-grid power fundamentals:
Calculating capacity and net revenue,''
\textit{J. Power Sources}, vol.~144, no.~1, pp.~268--279, Jun.~2005.
DOI:~10.1016/j.jpowsour.2004.12.025

\bibitem{v2g_barriers_2026}
M.~A.~Hannan \textit{et al.}, ``Vehicle-to-grid technology:
A review of the state of the art, challenges and opportunities,''
\textit{Renew. Sustain. Energy Rev.}, vol.~153, p.~111733, Jan.~2022.
DOI:~10.1016/j.rser.2021.111733

\bibitem{gao_economics_2025}
M. Gao \textit{et al.}, ``Economic viability of vehicle-to-grid
reassessed: A degradation cost integrated life-cycle analysis,''
\textit{Sustainability}, vol.~17, no.~12, p.~5626, Jun.~2025.
DOI:~10.3390/su17125626

\bibitem{liu_v2g_review_2026}
U.~ur~Rehman \textit{et al.}, ``Advancements and challenges in
vehicle-to-grid technology: A systematic literature review of energy
trading and optimization models and future directions,''
\textit{Energy Syst.} (Online First), Jan.~2026.
DOI:~10.1007/s12667-025-00772-1

\bibitem{haarnoja_sac_2018}
T. Haarnoja, A. Zhou, P. Abbeel, and S. Levine,
``Soft actor-critic: Off-policy maximum entropy deep reinforcement
learning with a stochastic actor,''
in \textit{Proc. 35th ICML}, Stockholm, Sweden, Jul.~2018,
pp.~1861--1870. arXiv:1801.01290

\bibitem{wang_sac_mg_2025}
Z.~Yu \textit{et al.},
``Energy optimization management of microgrid using improved soft
actor-critic algorithm,''
\textit{Int. J. Renew. Energy Dev.}, vol.~13, no.~2, pp.~329--339,
Mar.~2024.
DOI:~10.61435/ijred.2024.59988

\bibitem{hassan_drl_2026}
J.~Cao \textit{et al.},
``Deep reinforcement learning based energy storage arbitrage with accurate
lithium-ion battery degradation model,''
\textit{IEEE Trans. Smart Grid}, vol.~11, no.~5, pp.~4513--4521, Sep.~2020.
DOI:~10.1109/TSG.2020.2986333

\bibitem{chen_reward_2026}
Y.~Zhang \textit{et al.},
``Deep reinforcement learning-based battery conditioning hierarchical
vehicle-to-grid coordination for multi-stakeholder benefits,''
\textit{IEEE Trans. Smart Grid}, vol.~15, no.~2, pp.~1787--1800, Mar.~2024.
DOI:~10.1109/TSG.2023.3297474

\bibitem{wei_v2g_deg_2022}
Z. Wei \textit{et al.},
``Profitable vehicle-to-grid services with minimal battery degradation
using deep reinforcement learning,''
\textit{IEEE Trans. Smart Grid}, vol.~13, no.~5, pp.~3897--3911,
Sep.~2022.
DOI:~10.1109/TSG.2022.3168069

\bibitem{fang_deg_model_2024}
J. Fang \textit{et al.},
``A semi-empirical capacity fade model combining rainflow cycle counting
and Arrhenius calendar aging for NMC lithium-ion cells,''
\textit{J. Power Sources}, vol.~580, p.~233412, 2024.
DOI:~10.1016/j.jpowsour.2023.233412

\bibitem{smith_deg_freq_2026}
J.~Schmalstieg, S.~K\"{a}bitz, M.~Ecker, and D.~U.~Sauer,
``A holistic aging model for Li(NiMnCo)O$_2$ based 18650 lithium-ion
batteries,''
\textit{J. Power Sources}, vol.~257, pp.~325--334, Jul.~2014.
DOI:~10.1016/j.jpowsour.2014.02.012

\bibitem{osqp_2020}
B. Stellato \textit{et al.},
``OSQP: An operator splitting solver for quadratic programs,''
\textit{Math. Program. Comput.}, vol.~12, pp.~637--672, 2020.
DOI:~10.1007/s12532-020-00179-2

\bibitem{fan_marl_2024}
J.~Fan, H.~Wang, and A.~Liebman,
``Multi-agent reinforcement learning for decentralised electric vehicle
charging coordination with vehicle-to-vehicle energy exchange,''
in \textit{Proc. IEEE SmartGridComm}, Washington, DC, USA, Dec.~2023,
pp.~1--6.
DOI:~10.1109/SmartGridComm57891.2023.10333970

\bibitem{kempton_udel_v2g}
W. Kempton \textit{et al.},
\textit{V2G Power: Element of Energy Policy},
University of Delaware, Newark, DE, USA, Tech. Rep., 2001.


\bibitem{ieee_2030_8_2018}
IEEE Standard 2030.8-2018, \textit{Standard for the Testing of Microgrid
Controllers}, IEEE, New York, NY, USA, 2018.
DOI:~10.1109/IEEESTD.2018.8444947

\bibitem{typhoon_hil_2022}
Typhoon~HIL, \textit{Typhoon HIL\,604+ Real-Time Simulator: User Manual},
Typhoon HIL Inc., Somerville, MA, USA, 2022.
[Online]. Available: https://www.typhoon-hil.com

\bibitem{chil_v2g_2022}
A.~Zabetian-Hosseini, G.~Joos, and B.~Boulet,
``Controller hardware-in-the-loop validation of V2G control for DC fast
charging stations,''
in \textit{Proc. IEEE/PES T\&D Conf. Expo.}, New Orleans, LA, USA,
Apr.~2022, pp.~1--5.
DOI:~10.1109/TD43745.2022.9816830

\bibitem{opalrt_review_2024}
A.~Farrokhabadi \textit{et al.},
``Real-time simulation and hardware-in-the-loop testing based on
OPAL-RT ePHASORSIM: A review of recent advances and a simple validation
in EV charging management systems,''
\textit{Energies}, vol.~17, no.~19, p.~4893, Sep.~2024.
DOI:~10.3390/en17194893


\bibitem{nasa_battery_2007}
B.~Saha and K.~Goebel,
``Battery data set,''
NASA~Ames Prognostics Center of Excellence (PCoE), Moffett Field,
CA, USA, Dataset No.\,5-7 and 18, 2007.
[Online]. Available: https://ti.arc.nasa.gov/tech/dash/groups/pcoe/prognostic-data-repository/


\bibitem{lee_acndata_2019}
Z.~J.~Lee, D.~Chang, C.~Jin, G.~S.~Lee, R.~Lee, T.~Lee, and S.~Low,
``Large-scale adaptive electric vehicle charging,''
in \textit{Proc. IEEE Int. Conf. Commun., Control, Comput. Technol.
Smart Grids (SmartGridComm)}, Beijing, China, Oct.~2019, pp.~1--7.
DOI:~10.1109/SmartGridComm.2019.8909798


\bibitem{cohen_power_1988}
J.~Cohen,
\textit{Statistical Power Analysis for the Behavioral Sciences},
2nd~ed.\ Hillsdale, NJ, USA: Lawrence Erlbaum Associates, 1988.

\end{thebibliography}

\end{document}
